%% kapa: exam=1 round=20 year=2009 period=1 form=A subject=civil_law subject_ko=민법 seq=19 number=19
%% source: Q-Net 감정평가사 기출문제 (cst00309) artlSeq=5024178 file=extracted/1008/5024178/2009년 제20회 감정평가사 1차시험 기출문제/1교시.hwp
%% conversion: kordoc 4.24.51 -> tools/gen1.py (auto); review=proofread-sonnet; flags=none
\begin{question}{19}
대리에 관한 설명으로 옳지 \uline{않은} 것은?
\begin{choices}
  \item 대리행위의 하자는 대리인을 기준으로 판단하는 것이 원칙이다.
  \item 부동산등기의 신청에는 쌍방대리가 허용된다.
  \item 본인이 행위무능력자를 대리인으로 선임한 경우 본인은 대리인의 행위무능력을 이유로 대리행위를 취소할 수 없다.
  \item 대리인이 사망하면 복대리권은 소멸한다.
  \item 부동산매매계약 체결의 대리권을 가진 자는 특별한 사정이 없는 한 본인을 대리하여 계약을 해제할 수 있다.
\end{choices}
\end{question}
\answer{5}
